\section{Data and Valuation Setup}

\subsection{Controlled layer separation}
The integrated experiment is organized as a controlled handoff across three modules. Team 8 governs only the gold-price process; Team 4 supplies the operating forecasts; and Team 5 maps those inputs through a common tax, WACC, terminal-value and equity-bridge contract. In the four-model valuation comparison, only the gold engine changes. Barrick's Q1--Q2 2026 actuals and the same Team 4 production and unit-cost forecasts from Q3 2026 onward are frozen over 20 quarters; the same Team 5 corporate rules and WACC shocks are then reused across runs.



\subsection{Calibration snapshot and measure discipline}
The current calibration uses the 2 September 2026 GLD call surface acquired through the IB API. After the common DTE $\geq75$ and numerical filters, the dense surface contains 605 eligible observations across 12 expiries and 146 distinct strikes, with 79--653 days to expiry. The official $8\times8$ Chebyshev--Chebyshev sample contains 64 distinct actual contracts spanning eight expiries and 20 strikes. The risk-free input is a same-date NSS fit to 14 U.S. Treasury tenors \citep{nss1994,treasury2026}. Its 2.0608 bp RMSE is computed against the continuously compounded par-yield proxy; the curve is not described as a bootstrapped zero curve. The Barrick market reference is the 2 September close of USD~44.13, while the common gold anchor remains Barrick's Q2 2026 average realized price of USD~4,417/oz \citep{barrick2026}. The operating calendar is a fixed scenario index, not a forecast commencing on 2 September; earlier actuals are not future observations.

The option layer is estimated under the risk-neutral measure $\mathbb Q$. In the corporate experiment, only the option-implied distributional shape is transferred to the realized-price anchor. No GLD-share-to-ounce conversion or validated $\mathbb Q\rightarrow\mathbb P$ transformation is claimed. Combining this law with WACC defines a sensitivity operator, not risk-consistent pricing.

\begin{table}[H]
\centering
\caption{Frozen contract used for model attribution.}
\label{tab:frozen_contract}
\begin{tabularx}{\columnwidth}{@{}>{\bfseries}p{0.25\columnwidth}X@{}}
\toprule
Block & Current experiment \\
\midrule
Gold layer & Black--Scholes/GBM, Heston, Bates--Poisson and Full Bates--Hawkes calibrated by Team 8 under $\mathbb Q$. \\
Operations & Barrick Q1--Q2 2026 actuals followed by Team 4 deterministic production and unit-cost vectors from Q3 2026. \\
Corporate mapping & Team 5 tax, growth--ROIC--reinvestment schedule, stochastic WACC, terminal-value and equity-bridge logic. \\
Simulation & 8,192 paths, 260 fine steps, 20 quarters and a common WACC shock array. \\
\bottomrule
\end{tabularx}
\end{table}

\subsection{Option exercise convention}
GLD contracts are American-style \citep{cboe2026}. We set payout yield $q=0$: trust expenses reduce backing, not cash distributions \citep{spdr2025}. For non-distributing calls and positive rates, early exercise is suboptimal in the frictionless model. An 800-step CRR check on all 64 contracts gives an American--European premium below $10^{-8}$ USD; this tests that convention, not market frictions.

\subsection{Operating-model status}
Team 4's original operating study is richer than the object ultimately passed downstream. Cost of Sales is forecast from a log-space master chain; production is constructed bottom-up from ore processed, average grade and recovery. The unified valuation refreshes the handoff with 2026 actuals but still freezes the resulting 20-quarter production and cost vectors. Consequently, mine-level operating uncertainty is not propagated jointly with gold-price uncertainty. Team 4's legacy price simulation and illustrative EBITDA/valuation output are explicitly excluded.
